
\subsection{Limitations and Conclusion}

% Our method depends on accurate 4D geometry reconstruction. The OmniRe gaussian reconstruction backbone cannot perfectly handle the rolling shutter distortion, which will introduce some error of the geometry reconstruction especially when car drive very fast, but our diffusion model shows the strong robotness against some geometry errors. 
% Furthermore, our method is limited to editing the appearance of the scene not the geometry of the scene. Thus, future work can be improving the native 3DGS rasterizer to handle the rolling shutter issue and enable the geometry editing ability. 

Our method relies on 4D geometry reconstruction from OmniRe. While the Gaussian reconstruction backbone may introduce minor geometric inaccuracies because of rolling shutter distortion, particularly during high-speed driving, our diffusion model demonstrates strong robustness to such imperfections and still produces high-quality results. Additionally, our current framework focuses on appearance editing rather than geometry editing. Future work could explore integrating rolling shutter correction into the 3DGS rasterizer and extending the framework to support geometry editing capabilities.
% When the reconstructed geometry is imperfect, the synthesized results may degrade in quality. As shown in ~\Cref{fig:comparison}, inaccuracies in geometry can lead to visible artifacts, such as distorted lane markings or twisted effects on buildings.
% We presented GA-Drive, a novel framework for photorealistic driving simulation that enables novel trajectory synthesis and appearance editing through text prompts. By decoupling appearance from geometry, GA-Drive overcomes key limitations of existing reconstruction-based and diffusion-based simulators, including blurring results, lack of editability, and inconsistency with original views. Our enhanced 4D geometry reconstruction leverages both LiDAR and monocular estimated depth to achieve more accurate geometry reconstruction. We further introduced a dedicated pseudo-view synthesis pipeline that enables independent training of a video diffusion model, facilitating realistic and temporally coherent appearance generation under diverse weather conditions. To support long-term and large-shift camera trajectories, we proposed a Progressive Trajectory Expansion Strategy that incrementally extends novel view synthesis while maintaining consistency. Together, these innovations establish GA-Drive as a powerful tool for training and evaluating autonomous driving systems in diverse, photorealistic, and interactive environments. We believe that our framework opens new avenues for developing more robust and generalizable driving agents, especially under rare or challenging environmental conditions.


% Our method depends on accurate 4D geometry reconstruction. As shown in Fig.~\ref{fig:comparison}, inaccuracies in geometry can degrade synthesis quality, producing artifacts such as distorted lane markings or warped buildings. Additionally, appearance editing across frames may result in temporal incoherence if edits are inconsistent. Addressing these limitations is an important direction for future work.

We present GA-Drive, a novel framework for photorealistic driving simulation that enables novel trajectory synthesis and appearance editing. By decoupling appearance from geometry, GA-Drive overcomes key limitations of previous simulators, such as blurring, limited editability, and inconsistency with original views. Our improved 4D reconstruction yields more accurate geometry. A dedicated pseudo-view simulation pipeline enables training a video diffusion model without costly 4D reconstruction, allowing photorealistic novel view generation across diverse driving scenarios. To enable long-term novel view synthesis, we introduce a segment-wise video generation framework that ensures coherence across segments. Collectively, these innovations establish GA-Drive as a powerful tool for training and evaluating autonomous driving systems in diverse, photorealistic, and interactive environments.



